\chapter{Mathematics, Symbols, and Units}
\label{chap:math}

\begin{learningoutcomes}
  \item Type inline and displayed mathematics.
  \item Align multi-line calculations.
  \item Use consistent scientific units.
\end{learningoutcomes}

\section{Use math mode for mathematical meaning}

\begin{latexcode}{Complete example: essential mathematics}
\documentclass{article}
\usepackage{amsmath}
\usepackage{siunitx}

\begin{document}
Inline mathematics fits within a sentence: \(y = mx + c\).

\begin{equation}
  \bar{x} = \frac{1}{n}\sum_{i=1}^{n} x_i
  \label{eq:mean}
\end{equation}

The mean in Equation~\ref{eq:mean} uses \(n\) observations.

\begin{equation}
  E = mc^2
\end{equation}

The sample mass was \qty{12.5}{\milli\gram}, the temperature
was \qty{25}{\degreeCelsius}, and the interval was
\qtyrange{5}{10}{\second}.
\end{document}
\end{latexcode}

Use \code{\textbackslash(...\textbackslash)} for inline mathematics and a
numbered environment when the display will be discussed later.

\section{Common constructions}

\begin{latexcode}{Math construction recipes}
% Subscript and superscript
\(x_i^2\)

% Fraction, root, and Greek letters
\(\frac{a}{b}\), \(\sqrt{x}\), \(\alpha + \beta\)

% Sum, integral, and limit
\[
  \sum_{i=1}^{n} x_i
  \qquad
  \int_0^1 f(x)\,dx
  \qquad
  \lim_{x\to 0} f(x)
\]

% Text inside mathematics
\[
  y =
  \begin{cases}
    1, & \text{if } x > 0,\\
    0, & \text{otherwise.}
  \end{cases}
\]
\end{latexcode}

Braces group multi-character subscripts and powers:
\code{x\_\{max\}} and \code{e\^{}\{-kt\}}.

\section{Align related lines}

\begin{latexcode}{Aligned equations}
\begin{align}
  y &= mx + c, \\
  y - c &= mx, \\
  x &= \frac{y-c}{m}.
\end{align}
\end{latexcode}

The ampersand marks the alignment point. Use \code{align*} when numbers are not
needed.

\section{Vectors and matrices}

\begin{latexcode}{Vectors and matrices}
\[
  \mathbf{x} =
  \begin{bmatrix}
    x_1 \\
    x_2 \\
    x_3
  \end{bmatrix},
  \qquad
  A =
  \begin{pmatrix}
    1 & 0 \\
    0 & 1
  \end{pmatrix}.
\]
\end{latexcode}

\section{Units that remain consistent}

\begin{latexcode}{siunitx recipes}
\num{12345.678}
\si{\micro\metre}
\qty{9.81}{\metre\per\second\squared}
\qty{2.5e-3}{\mole\per\litre}
\qtyrange{18}{22}{\degreeCelsius}
\end{latexcode}

Keep descriptive words outside math mode. Let \pkg{siunitx} control spacing and
unit typography.

\begin{commonerror}
\textbf{Symptom:} a subscript, caret, or Greek command causes an error in prose.

\textbf{Fix:} place the expression in math mode, for example
\code{\textbackslash(\textbackslash alpha\_1\textbackslash)}.
\end{commonerror}

\begin{codelab}
Compile the essential-mathematics example. Add one aligned calculation, one
matrix, and three quantities from your subject. Refer to one numbered equation.
\end{codelab}

\begin{chapterchecklist}
  \item Variables are in math mode; ordinary words are not.
  \item Related equations align at meaningful symbols.
  \item Units are typed semantically with \pkg{siunitx}.
  \item Only equations discussed later are numbered.
\end{chapterchecklist}
